\section{Introduction}\label{sec:intro}

Kubernetes has become the de-facto platform for orchestrating containerised applications, and organisations increasingly rely on it to deploy, scale and operate production software \citep{rahman2023k8smisconfig,shamim2020xicommandments}. The behaviour of a Kubernetes workload is not defined in application code but in declarative configuration files, known as \emph{manifests}, which specify, among other things, how many replicas run, how the health of containers is probed, how much CPU and memory each container may consume, which privileges it holds and how it is exposed on the network. Because these settings directly determine the security posture and the availability of the running system, errors in manifests have operational and security consequences. Configuration mistakes can grant containers root or host-level privileges, leak secrets, allow resource exhaustion, or leave a service unavailable when a pod fails \citep{rahman2023k8smisconfig,zhang2026k8sconfigdefects,nsacisa2022k8shardening}. Writing correct, secure and reliable manifests is therefore an important but error-prone engineering task.

Empirical studies show that human-written manifests frequently contain such problems. Rahman et al.\ identified 1,051 security misconfigurations across 2,039 manifests mined from 92 open-source repositories and grouped them into 11 categories, such as absent resource limits, absent security contexts and privileged containers \citep{rahman2023k8smisconfig}. Zhang et al.\ later catalogued 15 categories of Kubernetes configuration defects, including probing, resource and scheduling defects, and found that eight widely used static analysis tools detect only 8 of these 15 categories \citep{zhang2026k8sconfigdefects}. In parallel, large language models (LLMs) such as ChatGPT, Claude and Gemini, together with open-weight models such as Llama, Qwen and DeepSeek, have become everyday assistants for software and infrastructure engineering. They are now routinely asked to produce Kubernetes manifests directly from natural-language descriptions \citep{zhang2024genaismells}. Research on general-purpose code has repeatedly shown that LLM assistants reproduce insecure patterns: around 40\% of programs generated by GitHub Copilot in security-relevant scenarios were vulnerable \citep{pearce2022asleep}, and around half of the functionally correct backends generated by state-of-the-art models were exploitable \citep{vero2025baxbench}. For infrastructure-as-code, syntactic validity and deployability have been found to be largely independent of security compliance \citep{santosvargas2026secureterraform,zhang2026iacgen}.

Despite this growing reliance on LLMs, the quality of LLM-generated Kubernetes configurations is still poorly understood. The first study on the topic analysed 98 ChatGPT-generated manifests collected from developer conversations and found that 35.8\% contained unset CPU and memory requirements \citep{zhang2024genaismells}. However, it relied on uncontrolled prompts, a single model from 2023, a single static analysis tool and no human baseline. A subsequent comparison with human-written manifests \citep{jcomsec2026aivshumank8s} re-used those 2023 ChatGPT measurements against an unrelated human corpus and, as its authors state, did not generate new manifests with newer models, relied on a single tool, and did not execute any manifest in a cluster. Benchmarks for LLM-generated cloud configuration, such as CloudEval-YAML \citep{xu2024cloudevalyaml} and FDE-Bench \citep{ding2026fdebench}, evaluate curated tasks rather than real projects and do not compare against the manifests that developers actually ship. Other Kubernetes-specific LLM research uses models to \emph{detect}, \emph{repair} or \emph{harden} existing manifests rather than to author them \citep{malul2024genkubesec,ye2025llmsecconfig,sgancohen2025kubeguard}. Consequently, there is no evidence on (i) how configurations produced by current LLMs compare with human-written configurations \emph{for the same applications}, (ii) how different LLMs compare with each other under identical requirements, and (iii) whether these configurations actually deploy and recover from failures in a real cluster, which is where their reliability ultimately matters.

The main objective of this study is to empirically evaluate the reliability and security of Kubernetes configurations generated by LLMs, using the configurations that the same real-world projects actually ship as a reference point. Specifically, the study addresses the following research questions:
\begin{itemize}
    \item \textbf{RQ1:} How do LLM-generated Kubernetes configurations compare with the human-written configurations shipped by the same projects in correctness, security, reliability, requirement satisfaction and runtime behaviour?
    \item \textbf{RQ2:} Which errors and misconfigurations occur most frequently in LLM-generated configurations?
    \item \textbf{RQ3:} How do different LLMs differ in the type and frequency of the problems they introduce, and how consistent is each model across repeated generations?
    \item \textbf{RQ4:} How reliably do LLM-generated configurations deploy and recover from failures in a real Kubernetes cluster compared with human-written configurations, and do static analysis findings predict these runtime outcomes?
    \item \textbf{RQ5:} Does explicitly requesting production-ready output change the security and reliability of LLM-generated configurations?
\end{itemize}

To answer these questions, this study selects seven non-trivial open-source applications whose maintainers ship Kubernetes manifests, ranging from stateful database-backed applications and multi-service microservice systems to a cluster add-on. For each application, every LLM receives the same \emph{docs bundle}: the project's own README, Dockerfile and Docker Compose file, copied verbatim from the repository at a pinned commit, with Kubernetes-specific passages (embedded manifests and deployment instructions) redacted. No specification is authored by the research team; the model sees exactly what the project's maintainers wrote, so the comparison cannot be confounded by a researcher's or an LLM's own wording. The same bundle is given to five LLMs: three proprietary models (OpenAI GPT, Anthropic Claude and Google Gemini) and two open-weight models (Qwen and DeepSeek). Because the human developers did not necessarily aim for production-grade manifests, two prompt variants are used. A neutral prompt asks only for manifests that deploy the application and is used for the comparison with the human-written configurations, which are additionally classified as production- or demonstration-oriented. A second prompt explicitly requests production-ready manifests and is used to measure the effect of the request itself. Each model generates three independent samples per project and prompt variant, and the resulting LLM-generated configurations are evaluated together with the human-written configurations shipped by the same projects. All configurations pass through the same automated evaluation pipeline, organised as a six-dimension benchmark. The dimensions are (D1) correctness, using kubeconform and server-side validation; (D2) security, using Kubescape, Trivy and the Kubernetes Pod Security Standards; (D3) static reliability, using kube-score, Polaris, KubeLinter and a reliability-controls checklist; (D4) requirement satisfaction against a fact sheet of the application's documented requirements; (D5) actual deployment on a disposable Kubernetes cluster with a functional smoke test; and (D6) resilience, measured through pod-deletion and container-crash fault-injection experiments that record recovery time and service availability. Results are compared across sources using non-parametric statistical tests and effect sizes.

The main contributions of this paper are as follows:
\begin{itemize}
    \item A \emph{project-matched} empirical comparison of human-written and LLM-generated Kubernetes configurations, in which every LLM receives the same real application's own documentation, with no researcher- or model-authored specification in between.
    \item A systematic comparison of five current proprietary and open-weight LLMs on Kubernetes configuration generation, including their consistency across repeated samples and their sensitivity to an explicit request for production-ready output.
    \item A six-dimension evaluation benchmark that links static correctness, security and reliability analysis to requirement satisfaction, real deployment and fault-injection-based recovery, making it possible to test whether static findings predict runtime behaviour.
    \item A replication package containing the docs bundles, prompts, raw model outputs, generated manifests, tool outputs and analysis scripts.
\end{itemize}

The rest of this paper is organised as follows. Section~\ref{sec:literature} reviews related work on Kubernetes misconfigurations, the security of LLM-generated code, LLM-based generation of infrastructure configurations, and LLM-based analysis and repair of Kubernetes manifests, and identifies the research gap. Section~3 describes the dataset, the generation procedure and the evaluation benchmark. Section~4 presents the experimental results for each research question. Section~5 discusses the findings and their implications for practitioners, and Section~6 discusses threats to validity. Finally, Section~7 concludes the paper and outlines directions for future research.
